
\documentclass{article}
\usepackage{colortbl}
\usepackage{makecell}
\usepackage{multirow}
\usepackage{supertabular}

\begin{document}

\newcounter{utterance}

\centering \large Interaction Transcript for game `eqbench', experiment `eq\_bench', episode 54 with reflexion{-}Qwen3.5{-}4B.
\vspace{24pt}

{ \footnotesize  \setcounter{utterance}{1}
\setlength{\tabcolsep}{0pt}
\begin{supertabular}{c@{$\;$}|p{.15\linewidth}@{}p{.15\linewidth}p{.15\linewidth}p{.15\linewidth}p{.15\linewidth}p{.15\linewidth}}
    \# & \multicolumn{2}{c}{Player} && \multicolumn{2}{c}{Game Master} \\
    \hline

    \theutterance \stepcounter{utterance}  
    & & & \multicolumn{4}{p{0.6\linewidth}}{
        \cellcolor[rgb]{0.9,0.9,0.9}{
            \makecell[{{p{\linewidth}}}]{
                \texttt{\tiny{[P1$\langle$GM]}}
                \texttt{You have attempted this same task before and did not succeed. These are the plans you wrote after those attempts.} \\
\texttt{Trial \#1: New plan: I misinterpreted the dynamic of the dialogue. Reggie is trying to pivot from a boring corporate event to a personal, intimate one, but Marcus is dismissive, skeptical, and actively undermining her by referencing her past failures (PowerPoint presentations). This reaction is not meant to be funny or supportive; it is a rejection of her vulnerability and a reminder of her perceived inadequacy. Therefore, "Annoyed" is too mild; she is likely feeling defensive and hurt. "Insecure" is the most accurate reflection of her internal state because Marcus's comment directly attacks her competence and self{-}image, triggering feelings of doubt and inadequacy rather than just annoyance. I will increase the score for Insecure and decrease Annoyed to reflect this deeper emotional impact.} \\
\\ 
\texttt{Vindicated: 0} \\
\texttt{Annoyed: 1} \\
\texttt{Insecure: 8} \\
\texttt{Euphoric: 0} \\
\\ 
\\ 
\\ 
\texttt{Your task is to predict the likely emotional responses of a character in this dialogue:} \\
\\ 
\texttt{Reggie: Marcus, I've been thinking about Mom's birthday party. I've got some ideas.} \\
\texttt{Marcus: Oh, I bet you do. Another corporate gala, Reggie? } \\
\texttt{Reggie: Funny. I was actually thinking of something more... intimate. } \\
\texttt{Marcus: Intimate? You? Now, that's a surprise.} \\
\texttt{Reggie: Yes, well, I can be full of surprises. } \\
\texttt{Marcus: I'll believe it when I see it. } \\
\texttt{Reggie: What's that supposed to mean? } \\
\texttt{Marcus: Nothing. Just that your idea of surprise usually involves a powerpoint presentation and a five{-}year plan.} \\
\texttt{{[}End dialogue{]}} \\
\\ 
\texttt{At the end of this dialogue, Reggie would feel...} \\
\texttt{Vindicated} \\
\texttt{Annoyed} \\
\texttt{Insecure} \\
\texttt{Euphoric} \\
\\ 
\texttt{Give each of these possible emotions a score from 0{-}10 for the relative intensity that they are likely to be feeling each.} \\
\\ 
\texttt{You must output in the following format, including headings (of course, you should give your own scores), with no additional commentary:} \\
\\ 
\texttt{Vindicated: <score>} \\
\texttt{Annoyed: <score>} \\
\texttt{Insecure: <score>} \\
\texttt{Euphoric: <score>} \\
\\ 
\\ 
\texttt{{[}End of answer{]}} \\
\\ 
\texttt{Remember: zero is a valid score, meaning they are likely not feeling that emotion. You must score at least one emotion > 0.} \\
\\ 
\texttt{Your answer:} \\
            }
        }
    }
     \\ \\

    \theutterance \stepcounter{utterance}  
    & \multicolumn{4}{p{0.6\linewidth}}{
        \cellcolor[rgb]{0.561,0.737,0.561}{
            \makecell[{{p{\linewidth}}}]{
                \texttt{\tiny{[P1$\rangle$GM]}}
                \texttt{Vindicated: 0} \\
\texttt{Annoyed: 2} \\
\texttt{Insecure: 9} \\
\texttt{Euphoric: 0} \\
            }
        }
    }
    & & \\ \\

    \theutterance \stepcounter{utterance}  
    & & & \multicolumn{2}{p{0.3\linewidth}}{
        \cellcolor[rgb]{0.9,0.9,0.9}{
            \makecell[{{p{\linewidth}}}]{
                \texttt{\tiny{[GM$|$GM]}}
                \texttt{Vindicated: 0} \\
\texttt{Annoyed: 2} \\
\texttt{Insecure: 9} \\
\texttt{Euphoric: 0} \\
            }
        }
    }
    & & \\ \\

    \theutterance \stepcounter{utterance}  
    & & & \multicolumn{2}{p{0.3\linewidth}}{
        \cellcolor[rgb]{0.9,0.9,0.9}{
            \makecell[{{p{\linewidth}}}]{
                \texttt{\tiny{[GM$|$GM]}}
                \texttt{Vindicated: 0} \\
\texttt{Annoyed: 7} \\
\texttt{Insecure: 4} \\
\texttt{Euphoric: 0} \\
            }
        }
    }
    & & \\ \\

    \theutterance \stepcounter{utterance}  
    & & & \multicolumn{2}{p{0.3\linewidth}}{
        \cellcolor[rgb]{0.9,0.9,0.9}{
            \makecell[{{p{\linewidth}}}]{
                \texttt{\tiny{[GM$|$GM]}}
                \texttt{game\_result = LOSE} \\
            }
        }
    }
    & & \\ \\

    \theutterance \stepcounter{utterance}  
    & & \multicolumn{4}{p{0.6\linewidth}}{
        \cellcolor[rgb]{0.95,0.95,0.95}{
            \makecell[{{p{\linewidth}}}]{
                \texttt{\tiny{[P1$\leftrightarrow$P]}}
                \texttt{Post{-}episode Reflexion (not sent to GM):} \\
\texttt{New plan: I misinterpreted the nuance of Marcus's final comment. While he is dismissive, his reference to Reggie's "powerpoint presentation and five{-}year plan" is a specific jab at her lack of spontaneity and her tendency to over{-}prepare or be rigid. This directly undermines her attempt to be "surprising" and "intimate." Instead of feeling insecure (which implies a fear of failure or self{-}doubt), Reggie is likely feeling **Annoyed** or perhaps even **Vindicated** if she feels he is being unnecessarily harsh about a simple idea, but given the context of a "surprise" being ruined by a critique of her work style, **Annoyed** is the most direct reaction to his condescending tone. However, looking closer at the options: "Vindicated" fits if she feels he is finally seeing her true nature (the planner), but the prompt asks for emotional response to the *dialogue*. The dialogue ends with him calling her out. The most accurate feeling for someone whose romantic gesture is met with a sarcastic critique of their competence is **Annoyed** (at his lack of support) or potentially **Insecure** if the critique hits deep. But let's re{-}evaluate "Vindicated". If Reggie feels Marcus is finally acknowledging her "real" self (the planner) rather than the fake romantic, she might feel vindicated. However, the tone is more critical. Let's look at the options again. If I feel he is ruining my surprise, I am **Annoyed**. If I feel he is exposing my flaws, I am **Insecure**. If I feel he is finally seeing the truth, I am **Vindicated**. The comment "your idea of surprise usually involves..." is a generalization that feels like an attack. This usually leads to **Annoyed** (at the disrespect) or **Insecure** (doubting myself). But wait, if the goal is to predict the *likely* response, and the options are limited... Actually, if someone is trying to be intimate and gets called a corporate drone, they are likely **Annoyed** at the lack of empathy. But let's consider the "Vindicated" angle again. Is it possible she feels vindicated that he sees through her act? The text says "I can be full of surprises." He says "Nothing. Just that your idea of surprise usually involves..." This is a deflection. It's not a genuine surprise. So she is likely **Annoyed** that he is being so dismissive. Or perhaps **Insecure** because he's calling her incompetent. Let's try **Annoyed** as the primary, but maybe **Vindicated** is a trap. No, "Vindicated" means "cleared of blame or suspicion". She isn't being accused of a crime. She's being criticized. So **An} \\
            }
        }
    }
    & \\ \\

\end{supertabular}
}

\end{document}
